\section*{Executive Summary}

This project reflects the ongoing development of the Digital Audio Workstation (DAW),
and it's built with C++, inspired by many popular DAWs like FL Studio, as well as
General MIDI. The goal of this project is to utilize a multi-threaded architecture that allows a group of
tracks to play in parallel with the other threads before being sent to the Master audio
effects and a mixer, optimizing perceptual audio quality while maintaining real-time constraints and synchronous timing under 
continuous workload.

The tracks include musical instruments equipped with Fluidsynth and Sfizz API,
allowing support for SoundFonts (SF2) stashed in Musical Artifacts, as well as SFZ, which are relatively
inexpensive to run and maintain. It currently has support for oscillator-based synthesizer with unison, which has
can offer a variety of sounds in fields of electronic music. Various IIR-based audio effects, such as
Biquad filters, State Variable Filters (SVF), Reverberation, compressor are also implemented for maintaining
professional sound quality, and are heavily based on Infinite Impulse Response (IIR) mechanism.

The test cases include perceptual listening, signal analysis, and real-time performance evaluation. 
All the Bode Plots and time-domain plots demonstrated the expected behaviour of the audio effects implemented. 
The development of the project, however, is currently underway, and is only limited to testing specified components.
So the theoretical estimation based on the multi-threaded DAG architecture are calculated, after 
measuring the time for each component. This still provided evidences that the mean and worst-case execution time
are well below the buffer time frame, even though it does not represent the actual stress test of a complete DAW 
operating with a full live production session, which will be done in the future.
